% TOP_PROBLEM: {"id": 3094, "title": "Edge-Unfolding Convex Polyhedra", "category_id": 6, "category": {"id": 6, "name": "geometry", "display_name": "Geometry", "description": "Euclidean and non-Euclidean geometry, geometric structures.", "slug": "geometry", "order_index": 6, "created_at": "2026-07-31T15:26:25.671Z"}, "status": "open", "problem_number": "OPG-60037", "set_id": 12}
% FIRST_POSTED: 2026-10-01
% ATTEMPT_STATUS: unresolved
% COMPLETION: 5
% MODEL: GPT 6 Astra Ultra
% UnsolvedMath ID 3094; source catalogue code OPG-60037
% !TeX program = lualatex
% Research attempt, 2026-10-01. Canonical statement preserved verbatim.
\documentclass[11pt,a4paper]{article}
\usepackage[margin=25mm]{geometry}
\usepackage{fontspec}
\setmainfont{lmroman10-regular.otf}[BoldFont=lmroman10-bold.otf,
 ItalicFont=lmroman10-italic.otf,BoldItalicFont=lmroman10-bolditalic.otf]
\usepackage{amsmath,amssymb,mathtools}
\usepackage{unicode-math}
\setmathfont{latinmodern-math.otf}
\AtBeginDocument{\let\setminus\smallsetminus}
\usepackage{xurl,hyperref}
\hypersetup{colorlinks=true,urlcolor=blue}
\setcounter{secnumdepth}{0}
\setlength{\parindent}{0pt}
\setlength{\parskip}{5pt}
\setlength{\emergencystretch}{3em}
\begin{document}
\section*{Edge-Unfolding Convex Polyhedra}\label{3094}
{\small UnsolvedMath ID 3094 · OPG-60037 · Geometry · OpenGarden}
\subsection{Definitions and mathematical statement}
Let $P\subset\mathbb R^3$ be a bounded convex polyhedron with nonempty
interior. Its boundary consists of finitely many planar polygonal faces,
meeting along its original edges and vertices. An edge-cut tree is a
spanning tree of the vertex-edge graph of $P$: it contains every vertex,
is connected, and has no cycle.

Cut the boundary surface along all edges of such a tree, duplicating the
cut edges on the resulting boundary. The remaining surface is a topological
disk made of the original polygonal faces, still attached along the uncut
edges. An edge unfolding places this disk into $\mathbb R^2$, isometrically
on each face, with the two incident faces of each uncut edge attached
along their common image. It can be obtained by successively rotating
faces about uncut edges until the surface is flat.

The conjecture states that every $P$ admits a choice of edge-cut tree
whose unfolded faces have pairwise disjoint interiors. This is the usual
nonoverlap condition for a net; contacts along face boundaries are allowed.
The flat result must be one connected piece containing every face exactly
once. Cutting through the interior of a face, adding new vertices to
supply extra cuts, or separating the faces into several independent pieces
is not allowed.

The problem is about existence of the final planar net. It does not
require that a physical motion from the original polyhedron to that net
avoid every intermediate collision. Convexity is imposed on the original
three-dimensional solid, not on the planar net, which may be nonconvex.
\subsection{Short English statement}
Can every convex polyhedron be cut only along an edge spanning tree and flattened into a single net with no overlapping face interiors?
\subsection{Research attempt}
\paragraph{Scope and literature check.}
Demaine's original posting [S2] and the primary nonconvex examples in [A1]
were checked on 1 October 2026. Nonconvex polyhedra can fail to edge-unfold,
even when individual faces are convex; those results do not settle this
convex-solid conjecture. The attempt below constructs a net for every right
prism over a convex polygon and audits the edge-tree requirement explicitly.
This is a familiar restricted family, not a new general theorem.

\paragraph{Construction of a connected net.}
Let $Q$ be a convex $m$-gon with successive edge lengths
$l_1,\ldots,l_m>0$, and let $P=Q\times[0,h]$, $h>0$. Assume no redundant
collinear vertices. Put $s_0=0$ and $s_i=\sum_{j\le i}l_j$. Develop the lateral
faces as the consecutive rectangles
\[
                       R_i=[s_{i-1},s_i]\times[0,h].
\]
They form one strip. Keep their consecutive common edges, but cut the final
seam between $R_m$ and $R_1$. Attach an isometric copy of the bottom polygon
along the edge $[0,l_1]\times\{0\}$, lying in $y\le0$. Attach the top polygon
along $[0,l_1]\times\{h\}$, lying in $y\ge h$.
Such placements exist because each attaching edge is a supporting edge of
the convex polygon. The copies may extend horizontally beyond the strip;
that is harmless. Their interiors lie in disjoint open half-planes, while
the strip interiors have $0<y<h$. Within the strip, rectangles have disjoint
interiors by their disjoint open horizontal intervals. Thus no two face
interiors overlap, and all faces remain in one connected piece.

\paragraph{The cuts are an original-edge spanning tree.}
There are $2m$ vertices and $3m$ original edges. Cut the one vertical seam,
and cut every top and bottom polygon edge except the two used for attachment
to $R_1$. The total number of cut edges is
\[
                       1+2(m-1)=2m-1.
\]
The remaining cut edges on each base form a path through all $m$ vertices,
because they are a polygon boundary cycle with one edge removed. The cut
vertical seam connects a vertex of one such path to its counterpart in the
other. Hence the cuts are connected on all $2m$ vertices and have no cycle:
they are exactly the required spanning tree. The uncut face-adjacency graph
is the path of $m$ rectangles with one base attached to each side of $R_1$;
it is also a tree. No face-interior cuts or added subdivision vertices occur.

\paragraph{A metric reason this does not prove the general case.}
The two decisive facts are that each lateral face is a rectangle of the
same height, and that unfolding consecutive faces places all their top edges
on $y=h$ and all bottom edges on $y=0$. These follow from the right-prism
geometry. They do not follow merely from a polyhedron having a belt of faces,
or from the uncut dual graph being a tree. For an oblique prism, even the
lateral parallelograms need not develop into a straight strip with two
parallel bounding lines. For a general convex polyhedron, there need not
be two distinguished base faces at all. A three-dimensional affine shear
also does not preserve the planar isometries used for unfolding; applying
such a shear to a successful spatial prism does not automatically shear its
net consistently across every face.

\paragraph{Verification and first remaining gap.}
At $m=3$ the construction has three rectangles, two triangles, and five cut
edges on six vertices; at $m=4$ it recovers a valid rectangular-box net with
seven cut edges on eight vertices. The proof checks interior separation by
half-planes rather than by a drawing, so extremely thin bases and large or
small heights introduce no exception. Boundary contacts are permitted by
the definition. For arbitrary convex polyhedra one still needs a choice of
edge tree together with a comparable global separation certificate for all
the developed faces. This attempt gives neither an algorithm that always
finds such a tree nor a counterexample.

\paragraph{Final status.}
\textbf{UNRESOLVED.} Right prisms over arbitrary convex polygons are handled
with an explicit net and edge-tree certificate; the full convex-polyhedron
assertion remains beyond the argument. Completion estimate: 5\%.

\subsection{Sources}
\begin{itemize}
\item[S1] ULAM.AI and the UnsolvedMath Contributors, supplied problems.json, record 3094 (OPG-60037), OpenGarden collection. Catalogue identity and source transcription; CC BY 4.0.
\url{https://huggingface.co/datasets/ulamai/UnsolvedMath}.
\item[S2] Open Problem Garden, Edge-Unfolding Convex Polyhedra. Original problem and discussion; the mathematical formulation below is expanded from the supplied source record.
\url{http://www.openproblemgarden.org/op/edge_unfolding_convex_polyhedra}.
\item[A1] M. Bern, E. D. Demaine, D. Eppstein, E. Kuo, A. Mantler and J. Snoeyink, \emph{Ununfoldable Polyhedra with Convex Faces}, primary author manuscript, abstract and introduction checked 1 October 2026. \url{https://erikdemaine.org/papers/Ununfoldable/paper.pdf}.
\end{itemize}
\end{document}
